% !TEX root = mixed_cubic_point_schemes.tex
\documentclass[a4paper,11pt]{ltjsarticle}
\usepackage[margin=25mm,headheight=16pt]{geometry}
\usepackage{amsmath,amssymb,amsthm,mathtools,bm,mathrsfs}
\usepackage{booktabs,longtable,array,enumitem}
\usepackage{xcolor}
\usepackage[unicode,colorlinks=true,linkcolor=blue!40!black,citecolor=blue!40!black,urlcolor=blue!45!black]{hyperref}
\usepackage{fancyhdr}
\usepackage{microtype}
\setlength{\parindent}{1em}
\setlength{\parskip}{0.25em}
\setlist{itemsep=0.2em,topsep=0.4em}
\allowdisplaybreaks[2]
\pagestyle{fancy}
\fancyhf{}
\fancyhead[L]{\small 混合三次超曲面の点スキーム}
\fancyhead[R]{\small 2026-11-13 / v1}
\fancyfoot[C]{\thepage}
\newcommand{\CC}{\mathbb C}
\newcommand{\ZZ}{\mathbb Z}
\newcommand{\PP}{\mathbb P}
\newcommand{\OO}{\mathcal O}
\newcommand{\rank}{\operatorname{rank}}
\theoremstyle{plain}
\newtheorem{theorem}{定理}[section]
\newtheorem{proposition}[theorem]{命題}
\newtheorem{lemma}[theorem]{補題}
\newtheorem{corollary}[theorem]{系}
\theoremstyle{definition}
\newtheorem{definition}[theorem]{定義}
\newtheorem{example}[theorem]{例}
\theoremstyle{remark}
\newtheorem{remark}[theorem]{注意}
\renewcommand{\proofname}{証明}

\hypersetup{pdftitle={混合三次非可換超曲面の点スキームとPoisson零点}}
\title{混合三次非可換超曲面の点スキームとPoisson零点}
\author{}
\date{2026年11月13日／研究草稿 v1}
\begin{document}
\maketitle
\begin{abstract}
中心一次元を二つ持つ五生成元の混合量子化族について、
中心三次商の点加群を表すスキームと、その半古典Poisson零点を比較する。
非零の量子化パラメータでは楕円曲線と長さ24の有限部分が現れ、
有限部分の非被約構造、射影空間内の接方向、混合係数に伴う衝突を記述する。
零量子化パラメータへの平坦閉包と、可換特殊ファイバー自体の点スキームは異なる。
中心三次元の閉公式とその自由代数内の証明を含め、
判別式上の六重点までを同じ原稿で扱う。
\end{abstract}

\section*{結果の構成と適用範囲}
扱う代数と中心元は第\ref{mt:sec:family}節で定義する。
Sklyanin楕円曲線の点加群と中心元の既知結果
\cite{PSPSmithTate,PSPDeLaet}を入力とし、
中心座標が非零のチャートは可換化によって直接計算する。
その非被約構造と埋め込みの相違を定理\ref{ps39:thm:comparison}で、
全点スキームを定理\ref{ps39:thm:global}で記述する。
二パラメータの平坦閉包は定理\ref{psl:thm:universal-flat}、
混合係数が零となる際の衝突は系\ref{psl:cor:collision}、
判別式上の局所環は定理\ref{pdb:thm:sixfold}にまとめる。

本稿は指定された混合族のスキーム論的な研究であり、
任意の非可換三次超曲面の分類や残余圏の可換復元を主張しない。
点スキームとPoisson零点の同一性を非零パラメータの各ファイバーで主張することと、
平坦閉包の特殊ファイバーを同定することも区別する。
初出・新規性の包括的な照合は未完了である。

\section{最小混合試験族の三次中心元の閉公式}
\label{mt:sec:family}

一楕円ブロックの分離型から一般混合型へ進むため、
五変数のcubic
\begin{equation}\label{mt:eq:classical}
 F_\lambda=\frac{x^3+y^3+z^3}{3}+u^3+v^3+\lambda uxy
\end{equation}
を固定する。$u,v$を中心、$h,\lambda$を次数零のパラメータとし、
\begin{equation}\label{mt:eq:relations}
 [y,z]=h(x^2+\lambda uy),\qquad
 [z,x]=h(y^2+\lambda ux),\qquad
 [x,y]=hz^2
\end{equation}
で周囲の代数$A_{h,\lambda}$を定める。
本節ではこの族の三次中心元を多項式の閉公式で構成し、
自由代数内の三つの恒等式によって中心性を証明する。
これは可換復元先の二次係数$G_2$の計算ではなく、
それを計算するための非可換代数側の入力である。
中心式の特殊形は既存の混合中心式の有限体検算に現れるが、
ここでは特殊化を用いず、全パラメータに対する証明書を与える。
閉公式の先行研究上の初出は未確認である。

\begin{lemma}[古典特殊ファイバーの滑らかさ]
\label{mt:lem:classical-smooth}
$V(F_\lambda)\subset\PP^4$は$\lambda^3\ne-3$のとき、かつそのときに限り滑らかである。
\end{lemma}
\begin{proof}
全偏微分が零なら$z=v=0$であり、残る方程式は
\[
 x^2+\lambda uy=y^2+\lambda ux=3u^2+\lambda xy=0.
\]
$u=0$なら$x=y=0$となるので射影点は存在しない。
$u=1$とすると、$\lambda=0$では解がなく、$\lambda\ne0$なら$x,y$は非零である。
最初の二式を掛けると$xy=\lambda^2$を得て、第三式は$3+\lambda^3=0$となる。
逆にこの等式が成立すれば、$x^3=3$を満たす$x$と
$y=-x^2/\lambda$を選ぶことで解を得る。
三次のEuler等式により、これらの解は$F_\lambda=0$も満たす。
\end{proof}

\subsection{中心三次式と自由代数の証明書}

$a=\lambda u$と略記し、$T=xyz+xzy+yxz+yzx+zxy+zyx$と置く。

\begin{proposition}[混合試験族の中心元]
\label{mt:prop:center}
\begin{equation}\label{mt:eq:center}
 \boxed{\begin{aligned}
 \Omega_{h,\lambda}={}&
 \frac{x^3+y^3+z^3}{3}-\frac{h^2}{9}T
 +\frac a2(xy+yx)\\
 &+\frac{h^2a^2}{6}z+\frac{h^4a}{18}z^2+u^3+v^3.
 \end{aligned}}
\end{equation}
は$A_{h,\lambda}$の三次中心元である。
これは$\CC[h,\lambda]$上の恒等的な中心性であり、
$h\ne0$への局所化も、有限$h$-jetへの制限も必要としない。
$h=0$で$\Omega_{0,\lambda}=F_\lambda$である。
\end{proposition}
\begin{proof}
まず$a$を独立な中心係数とし、自由代数
$\mathbb Q[h,a]\langle x,y,z\rangle$で計算する。
関係式を
\[
 r_1=xy-yx-hz^2,\quad
 r_2=xz-zx+hy^2+hax,\quad
 r_3=yz-zy-hx^2-hay
\]
と書く。二つの演算を
\[
 \mathcal S_p(r)=rp^2+prp+p^2r,\qquad
 \mathcal D_{p,q}(r)=rpq-prq-qrp+qpr
\]
とし、$\{p,q\}_+=pq+qp$とする。
$\omega=\Omega_{h,\lambda}-u^3-v^3$について、以下の三式が
\emph{関係式で割る前の自由代数}において成立する：
\begin{align}
 [\omega,x]={}&
 -\frac13\mathcal S_y(r_1)-\frac13\mathcal S_z(r_2)
 +\frac{h^2}{18}
   \bigl(\mathcal D_{z,x}(r_1)+\mathcal D_{y,x}(r_2)\bigr)\notag\\
 &+\frac{ah^3}{18}[x,r_1]-\frac a2\{x,r_1\}_+
  +\frac{ah}{6}[z,r_2]+
  \frac h6[yz+zy,r_3],\label{mt:eq:certificate-x}\\
 [\omega,y]={}&
 \frac13\mathcal S_x(r_1)-\frac13\mathcal S_z(r_3)
 -\frac{h^2}{9}\mathcal D_{z,y}(r_1)
 +\frac{ah^3}{18}[y,r_1]+\frac a2\{y,r_1\}_+\notag\\
 &+\frac{ah^2}{18}[r_2,x]+\frac h9[r_2,2xz+zx]
 +\frac{h^2}{9}[r_2,y]y+\frac{ah}{6}[r_3,z],
 \label{mt:eq:certificate-y}\\
 [\omega,z]={}&
 \frac{a(2h^3+9)}{18}[r_1,z]
 +\frac h9[2xy+yx,r_1]+\frac{h^2}{9}[r_1,z]z\notag\\
 &+ayr_2+\frac13\mathcal S_x(r_2)
 -\frac{h^2}{9}\mathcal D_{y,z}(r_2)
 +ar_3x+\frac13\mathcal S_y(r_3).
 \label{mt:eq:certificate-z}
\end{align}
各式は語と$h,a$の冪ごとに展開して確認でき、全係数は有理数である。
右辺の各項は$r_i$を含むので、商代数で全交換子が零となる。
$a=\lambda u$を代入し、中心スカラー$u^3+v^3$を戻せば中心性を得る。
零パラメータでの表示は直接代入による。
\end{proof}

\subsection{Artin切断でのsorted PBW表示}

閉公式の中心性と、周囲の代数のflatnessを別々に確認する。
$u,v$の中心性を先に使い、active語を$z<y<x$の順に並べる。
還元規則は
\begin{equation}\label{mt:eq:rewrite}
 xy\longmapsto yx+hz^2,\qquad
 xz\longmapsto zx-hy^2-hax,\qquad
 yz\longmapsto zy+hx^2+hay
\end{equation}
である。これは一般の非零$h$で通常の語順序に関して停止する規則とは主張しない。

\begin{lemma}[全Artin切断でのPBW性]
\label{mt:lem:artin-pbw}
任意の$N\ge0$について、$h^{N+1}=0$を課した代数は
$\CC[\lambda,u,v,h]/(h^{N+1})$上で
$\{z^iy^jx^k:i,j,k\ge0\}$を基底に持つ。
従って各次数成分は$\CC[\lambda,h]/(h^{N+1})$上自由であり、
周囲のHilbert級数は$(1-t)^{-5}$である。
\end{lemma}
\begin{proof}
$h^r w$の還元に、辞書式に並べた二つの非負整数
$(N-r,\operatorname{inv}(w))$を用いる。
$h$を含まない交換項では反転数が減り、補正項では$h$の冪が増える。
従って$h^{N+1}=0$のもとで全還元が停止する。
active語の重なりは$xyz$の一つだけである。
最初に$xy$を還元した場合と、最初に$yz$を還元した場合は、いずれも有限回の置換で
\begin{equation}\label{mt:eq:overlap}
 zyx+h(x^3-y^3+z^3)
\end{equation}
になる。$a$を含む項は両経路でそれぞれ相殺する。
包含の曖昧さはない。係数に逆元を導入しないmonicな合流的還元なので、
既約語$z^iy^jx^k$が基底となる。
中心係数$a$を$\lambda u$へ置き換え、独立な中心変数$v$を加えても同じ証明が成り立つ。
\end{proof}

この補題は全ての$h$-Artin切断およびその逆極限についてのformal PBW性を示す。
非零$h$での全ファイバーのAS-regular性や代数的族全体のflatnessを、
Artin切断の計算だけから結論していない。
また、中心商は$v$に関してmonicな三次式であるため、
各Artin切断で$1,v,v^2$を、$v$を含まない周囲の代数上の基底に持つ。
従って中心三次商も次数別に平坦で、Hilbert級数は
$(1-t^3)/(1-t)^5$である。



\paragraph{補足の比較式.}
以下で形式復元候補と呼ぶ三次式は、$\lambda=1$における
次の比較用多項式を指す：
\begin{equation}\label{ps:comparison-cubic}
\begin{split}
 \widetilde G_h=G_h+v^3={}&(x^3+y^3+z^3)/3+v^3+h^2xyz/3\\
 &+(1+h^6/54)uxy-h^4uz^2/18-h^2u^2z/12\\
 &+(1+29h^6/648+h^{12}/4374)u^3.
\end{split}
\end{equation}
本稿で用いるのは、この式の$u=v=0$制限が
$(x^3+y^3+z^3+h^2xyz)/3$に等しいことだけである。
この三次式との圏同値の存在は本稿の結論に含めない。

% 2026-11-11 / v38
\section{混合cubicのpoint schemeとPoisson零点集合}
\label{sec:points-poisson-cubic}

非可換のpoint schemeは、一般にはPoisson構造の零点集合と同じではない。
ここでは固定パラメータでの集合論的な記述を与え、h=0の極限、
中心三次商、形式復元先のcubicを区別する。

\subsection{混合中心三次商}
\[
F_1=(x^3+y^3+z^3)/3+u^3+v^3+uxy,
\]
\[
[y,z]=h(x^2+uy),\quad[z,x]=h(y^2+ux),\quad[x,y]=hz^2,
\quad[u,-]=[v,-]=0.
\]
この関係で定まる周囲の代数を$A_h$とし、
第\ref{mt:sec:family}節で構成した中心元$\Omega_h$による商を
$B_h=A_h/(\Omega_h)$とする。標数零の代数閉体において
$h\ne0$、$h^6\ne-27$とする。点モジュールはHilbert級数
$(1-t)^{-1}$の次数付き巡回右加群とし、標準の射影列$(p_i)_{i\ge0}$で表す。
ここでは集合論的なpoint parametrizationを扱い、
truncated point schemeの普遍族や非被約構造は同定しない。

\begin{proposition}[混合族のpoint moduleの台]\label{prop:psp-point-support}
上のパラメータでは、$B_h$の点モジュールは、その初期射影点によって、
次の二種類に分類される。
\[
\begin{split}
E_h&=\{u=v=0,\ x^3+y^3+z^3+h^2xyz=0\}
             \subset\mathbb P^4,\\
\Sigma&=\{z=0,\ x^2+uy=y^2+ux=0,\ F_1=0,\ u\ne0\}.
\end{split}
\]
$E_h$は滑らかなSklyanin楕円曲線で、$\Sigma$は12点からなる。
従ってpoint moduleの初期点の集合は
\[
|\operatorname{Point}(B_h)|=E_h\sqcup\Sigma.
\tag{PSP1}\label{eq:psp-support}
\]
$u$または$v$が非零の点モジュールはすべて定数列であり、
$E_h$上のものはSklyaninのtranslationで推移する。
\end{proposition}
\begin{proof}
各中心一次元$c=u,v$との交換関係を点モジュールの隣接座標
$p_i=(x_i,y_i,z_i,u_i,v_i)$、$p_{i+1}$へ多重線形化すると、
$c_ip_{i+1,a}=p_{i,a}c_{i+1}$（全一次生成元$a$）となる。
$c_i\ne0$なら$p_{i+1}=p_i$であり、
$c_{i+1}\ne0$も従う。前後の射影点にも同じ議論を適用するので、
中心座標がいずれかで非零の場合は全て定数列である。

定数列では交換子が零となるため、$h\ne0$から
\[
z^2=x^2+uy=y^2+ux=0
\]
を得る。点への中心元の作用は非可換語を同じスカラーに置く通常の代入である。
$z=0$を用いるとその値は$F_1$に一致する。
$u=0$なら$x=y=z=0$、$F_1=v^3\ne0$で不可能。
$u=1$に規格化すると、
$x=y=0$の場合は$v^3=-1$で三点。
$x,y\ne0$の場合は$xy=1$、$x^3=-1$、$y=-x^2$となり、
$v^3=-4/3$なので九点。いずれも定数列を定義する
（全関係式と中心元のスカラー代入が零）。

$u=v=0$の点列は三生成元のSklyanin代数
\[
S_h=\mathbb C\langle x,y,z\rangle/
(yz-zy-hx^2,zx-xz-hy^2,xy-yx-hz^2)
\]
の点列と同じである。この代数はSklyaninパラメータ
$(a,b,c)=(1,-1,-h)$を持つ。$h^6\ne-27$では
Sklyaninのpoint schemeが滑らかな楕円曲線
\[
abc(x^3+y^3+z^3)-(a^3+b^3+c^3)xyz=0
\]
であり、代入すると$E_h$を得る
\cite{PSPSmithTate,PSPDeLaet}.
この代数の中心三次元は全点モジュールを消す。
本研究の$\Omega_h\bmod(u,v)$はその非零スカラー倍なので、
すべての$S_h$点モジュールが$B_h$へ降りる。
中心座標が常に零の場合に他の点列はない。以上を合わせる。
\end{proof}

\begin{proposition}[Poisson零点との比較]\label{prop:psp-comparison}
$h=0$における$cubic$ $X_0=V(F_1)$のPoissonテンソル
\[
\{y,z\}=x^2+uy,\quad\{z,x\}=y^2+ux,\quad\{x,y\}=z^2
\]
の零点集合は台として
\[
|D_0(\pi_{X_0})|=E_0\sqcup\Sigma,
\quad E_0=\{u=v=0,\ x^3+y^3+z^3=0\}.
\tag{PSP2}\label{eq:psp-poisson}
\]
ここで臨界点部分の零点スキームは$24$の長さをもち、
$\Sigma$の各点で横断長2である。
一般の$h\ne0$では$E_h\ne E_0$であり、
$\operatorname{Point}(B_h)$は$X_0$に含まれさえしない。
一方、既存の有限対角から得た形式復元候補$G_h+v^3$について、
その$u=v=0$制限の零軌跡は$E_h$に等しい。
\end{proposition}
\begin{proof}
$E_0$で各活性JacobianのEulerベクトルが像に属するため、
射影化された双ベクトルは零になる。
$u=v=0$以外では中心方向が非零なので、射影化で階数は減らず、
Poisson零点の方程式は
\[
x^2+uy=y^2+ux=z^2=F_1=0.
\]
これは$\Sigma$の台を定める。
$u=1$上で$x^2+y=y^2+x=0$の解は
$(x,y)=(0,0)$または$(x^3,y)=(-1,-x^2)$で四個であり、
$x,y$方向のJacobian行列式はそれぞれ$-1$、$3$で非零。
各点で$v^3=-1$または$-4/3$なので$v$方向も単純である。
従って$z^2$が残る唯一の厚化で、12点それぞれ長さ2、合計24である。

$E_h$と$E_0$は$h\ne0$なら異なる三次曲線である。
例えば$h=1$では$p=[1:1:-1:0:0]$が$E_1$に属すが
$F_1(p)=1/3\ne0$。一般の$h$でも$xyz\ne0$の
$E_h$上の点は$F_1=-h^2xyz/3\ne0$。
最後に既存候補の$u=v=0$制限は
$(x^3+y^3+z^3+h^2xyz)/3$であるから一致する。
\end{proof}

\begin{remark}[埋め込みとスキーム論的分類]
$h=0$で$B_0$は可換cubic $X_0$の斉次座標環なので、
そのpoint schemeは$X_0$全体であり$D_0(\pi)$と等しくない。
$E_h\to E_0$という曲線の極限からpoint scheme全体の
flatな特殊化を結論してはならない。
本節の集合論的分類は次節でスキームとして精密化する。
孤立12点は各々長さ2の非被約点でPoisson零点の孤立厚化に
抽象的には同型だが、射影空間内の接方向は一般に異なる。
三次元Sklyaninのshiftは楕円曲線上の非自明なtranslationであり、
Poisson零点上の定数列とは異なる非可換動力学を保つ。
\end{remark}

% 2026-11-11 / v40
\section{混合cubicのpoint schemeの非被約構造}
\label{ps39:sec}
以下では混合族$\lambda=1$、$h\in\mathbb C^\times$、
$h^6\ne-27$を固定する。$B_h=A_h/(\Omega_h)$とし、
point schemeとはHilbert級数$(1-t)^{-1}$のgraded右point moduleの
初期点を表す射影スキームを意味する。
スキーム論的な計算と台集合の計算を区別する。

\subsection{中心一次元の非零チャート}
\begin{lemma}[中心座標による点列の対角化]\label{ps39:lem:abelian}
$u\ne0$におけるpoint moduleの族は、スキームとして
$B_h$の可換化の$u\ne0$チャートと一致する。
\end{lemma}
\begin{proof}
可換試験環$T$上の点列を$p_i=(x_i:y_i:z_i:u_i:v_i)$と表す。
多重線形化した中心性$[u,a]=0$（任意の一次生成元$a$）から
$u_ia_{i+1}-a_iu_{i+1}=0$を得る。$u_i$が単元なら各点を
$u_i=1$に規格化でき、$a_{i+1}=a_i$となる。
全$i$について繰り返すと点列はスキーム論的に一定である。
残る多重線形化関係式は、すべて同一座標への可換代入に等しい。
逆に可換化の任意の$T$値点は一定点列を定義する。
この対応は$T$に自然なので主張が従う。
\end{proof}

\begin{proposition}[孤立部分は12個の二重点]\label{ps39:prop:length}
$u=1$におけるpoint schemeのイデアルは
\begin{equation}\label{ps39:eq:ideal}
I_h=\left(x^2+y,\ y^2+x,\ z^2,\
1+v^3+\frac{xy}{3}
+h^2z\left(\frac16-\frac{2xy}{3}\right)\right).
\end{equation}
これは12個の相異なる点に支持を持つ長さ2の非被約点の非交和で、
総長は24である。
\end{proposition}
\begin{proof}
$h\ne0$なので交換子を可換化すると
$x^2+uy=y^2+ux=z^2=0$となる。
既存の混合族の中心三次元の可換化は
\[
\Omega_h^{\mathrm{ab}}
=F_1-\frac{2h^2}{3}xyz+\frac{h^2u^2}{6}z
+\frac{h^4u}{18}z^2.
\]
$u=1$で最初の関係を使うと$x^3=y^3=-xy$、$z^3=0$であり、
この式が\eqref{ps39:eq:ideal}へ変わる。
$x^2+y=y^2+x=0$は$(x,y)=(0,0)$、または
$x^3=-1,\ y=-x^2$という計四個の解を持つ。
Jacobian行列式は$4xy-1$で、前者で$-1$、後者で$3$だから
すべて被約である。$xy=0$または$1$に従い、
$z=0$で$v^3=-1$または$-4/3$となる。
いずれも三つの相異なる非零根を持ち、支持は計12点。
各点で$3v^2$が単元なので最後の式は$v$を$z$の冪零変形として
一意に解く。残る条件は$z^2=0$のみであり、
各局所環は$\mathbb C[z]/(z^2)$に同型である。
\end{proof}

\subsection{Poisson零点との同型と埋め込みの相違}
\begin{theorem}[孤立部分の比較]\label{ps39:thm:comparison}
$X_0=V(F_1)$上のPoisson零点の孤立部分を$Z_{\mathrm{crit}}$、
非可換point schemeの孤立部分を$\Sigma_h$とする。
このとき$\Sigma_h\simeq Z_{\mathrm{crit}}$がスキームとして成立するが、
$h\ne0$では両者は$\mathbb P^4$内の同じ閉部分スキームではない。
$u=1$での同型は
\begin{equation}\label{ps39:eq:shear}
(x,y,z,v)\longmapsto
\left(x,y,z,\
v'=v+\frac{h^2z}{3v^2}
 \left(\frac16-\frac{2xy}{3}\right)\right)
\end{equation}
で与えられる。
\end{theorem}
\begin{proof}
Poisson零点の孤立部分は
\[
I_0=(x^2+y,\ y^2+x,\ z^2,\ 1+v^3+xy/3)
\]
で定義される。両スキーム上で$v^3=-1$または$-4/3$を法とする
冪零元の補正だけがあるので、$v$は単元である。
$z^2=0$より
\[
(v')^3=v^3+h^2z(1/6-2xy/3).
\]
従って\eqref{ps39:eq:shear}で$I_h$と$I_0$が対応し、
逆写像は$z$の一次のせん断を逆向きに施すことで得られる。

一方、$(x,y,z)=(0,0,0)$、$v_0^3=-1$の点で
Poisson零点の埋め込み接方向は$\delta v=0$、
point schemeの接方向は
$3v_0^2\delta v+(h^2/6)\delta z=0$である。
$h\ne0$では異なるため、閉部分スキームとしては一致しない。
\end{proof}

\subsection{Sklyanin曲線と全point scheme}
\begin{theorem}[一般非零パラメータのpoint scheme]\label{ps39:thm:global}
$h\ne0$、$h^6\ne-27$のとき、point moduleの初期点を表す
スキームは
\[
\boxed{\quad
\operatorname{Point}(B_h)
\simeq E_h\sqcup\coprod_{a=1}^{12}
\operatorname{Spec}\mathbb C[\varepsilon]/(\varepsilon^2).
\quad}
\]
ここで$E_h=V(u,v,x^3+y^3+z^3+h^2xyz)$は滑らかな楕円曲線である。
特に$h=0$のpoint schemeは可換cubic $X_0$全体であって、
上の表示が$h=0$までpoint scheme全体の平坦族をなすとはしない。
\end{theorem}
\begin{proof}
$u\ne0$では直前の命題を使う。
$u=0,v\ne0$なら中心性から点列は一定であるが、
可換化した交換子は$x=y=z=0$を強制し、
中心三次元の値は$v^3\ne0$なのでpoint moduleは存在しない。
残りの支持は$u=v=0$であり、三生成元Sklyanin代数
\[
S_h=\mathbb C\langle x,y,z\rangle/
(yz-zy-hx^2,\ zx-xz-hy^2,\ xy-yx-hz^2)
\]
のpoint schemeとなる。既知のSklyanin定理
\cite{PSPSmithTate,PSPDeLaet}によりこれは滑らかな$E_h$で、
その中心三次元はすべてのpoint moduleを消す。
$\Omega_h\bmod(u,v)$はその中心三次元の非零倍なので、
これらのpoint moduleは$B_h$へ降りる。

さらに$E_h$方向に$u,v$の冪零的な厚化が生じないことを示す。
Sklyaninのshiftは$E_h$上の平行移動で、$h\ne0$なら恒等ではない。
従って$p_0$と$p_1$は$E_h$のどの点でも異なる。
その点の局所環である二次小行列
$p_{0,a}p_{1,b}-p_{0,b}p_{1,a}$が単元になる。
中心性の多重線形化
$u_0p_{1,a}=p_{0,a}u_1$および
$v_0p_{1,a}=p_{0,a}v_1$を二座標について消去すると、
その小行列と$u_0,v_0$との積は零となる。
従って局所環上で$u_0=v_0=0$が厳密に成立し、
近傍のpoint schemeはSklyaninの被約$E_h$に等しい。
$E_h$と$u\ne0$部分は互いに交わらないので非交和となる。
\end{proof}

\begin{remark}[結論の範囲]
Poisson零点全体$E_0\sqcup Z_{\mathrm{crit}}$と
非可換point scheme$E_h\sqcup\Sigma_h$は一般には一致しない。
孤立部分は抽象スキームとして同型だが、法線方向の埋め込みは異なる。
復元cubicの$u=v=0$制限は$E_h$に等しいものの、
point scheme全体が復元cubic上にあるとは主張しない。
Sklyanin退化値$h^6=-27$と$h=0$は上の定理の範囲外である。
\end{remark}

\subsection{零変形パラメータへの平坦極限}
\begin{theorem}[point schemeの平坦閉包はPoisson零点]
\label{ps39:thm:flat-limit}
$U=\operatorname{Spec}\mathbb C[h,(h^6+27)^{-1}]$とし、
$\mathcal Z\subset\mathbb P^4_U$を次の二つの閉部分スキームの非交和とする。
\[
\mathcal E=V(u,v,x^3+y^3+z^3+h^2xyz),
\]
\[
\mathcal S=V(x^2+uy,y^2+ux,z^2,\Omega_h^{\mathrm{ab}}),
\quad
\Omega_h^{\mathrm{ab}}=
F_1+h^2z\left(\frac{u^2}{6}-\frac{2xy}{3}\right)
+\frac{h^4u}{18}z^2.
\]
すると$\mathcal Z=\mathcal E\sqcup\mathcal S$は$U$上平坦で、
\[
\boxed{\quad
\mathcal Z_h\simeq\operatorname{Point}(B_h)\quad(h\ne0),
\qquad
\mathcal Z_0=D_0(\pi_{X_0}).
\quad}
\tag{PS40}\label{ps39:eq:flat}
\]
ここで第二式は$\mathbb P^4$の閉部分スキームとして成立する。
従って$h\ne0$のpoint schemeの$\mathbb P^4_U$内の
スキーム論的閉包の特殊ファイバーはPoisson零点スキームである。
全ファイバーのHilbert多項式は$3t+24$である。
一方$B_0$自体のpoint schemeは可換cubic $X_0$全体であり、
$\mathcal Z_0$とは異なる。
\end{theorem}
\begin{proof}
$\mathcal S$の$u=0$部分では三つの二次方程式から
$x=y=z=0$、中心三次式から$v^3=0$となる。
射影的共通零点は存在しないので
$\mathcal S\subset\{u\ne0\}$である。
従って$\mathcal E$と$\mathcal S$は交わらない。

$u=1$の$\mathcal S$は式\eqref{ps39:eq:ideal}で定義される。
この環で$v$は単元であり、
式\eqref{ps39:eq:shear}は$h$を固定した
$\mathbb C[h]$-代数の同型となる。
その逆は$z^2=0$の下で同じせん断を逆向きに施す式で与えられる。
従って$\mathcal S\simeq\mathcal S_0\times U$で、
長さ24の有限平坦族である。
$\mathcal E$は$\mathbb P^2_U$内の相対三次曲線で、
$h^6+27$を反転したため滑らかな楕円曲線族となり、特に平坦である。
よって非交和$\mathcal Z$も平坦である。

$h\ne0$では定理\ref{ps39:thm:global}を使い、
$\mathcal Z_h$がpoint schemeと一致する。
$h=0$では$\mathcal E_0=E_0$、
$\mathcal S_0=V(x^2+uy,y^2+ux,z^2,F_1)$となり、
前節の零点計算から
$\mathcal Z_0=E_0\sqcup Z_{\mathrm{crit}}=D_0(\pi_{X_0})$
がスキームとして成立する。
$U$は整域で$\mathcal Z$は$U$-flatなので、
$h=0$に支持される余分な閉成分はなく、
$\mathcal Z$は穿孔底上のpoint schemeのスキーム論的閉包である。
各$\mathcal E_h$は次数3・種数1の平面三次曲線だから
Hilbert多項式は$3t$、$\mathcal S_h$は長さ24だから、
全体の多項式は$3t+24$である。
$h=0$の非可換関係は可換化し、$B_0$のpoint schemeが
$X_0$全体となることも明らかである。
\end{proof}

% Flat point-scheme family with varying mixed coefficient.
\section{混合係数を動かしたpoint schemeの同時平坦極限}
\label{psl:sec:family}

前節の$\lambda=1$のpoint schemeの構成を、混合係数$\lambda$そのものを動かす二パラメータ族に拡張する。
特に$\lambda=0$では孤立点の支持が合体するが、scheme長は保存される。
中心三次商のpoint schemeと、その半古典Poisson構造の零点スキームの関係は、
この衝突を含む族でも維持される。

\subsection{二パラメータの有限平坦模型}

$\Lambda=\operatorname{Spec}\mathbb C[\lambda,(\lambda^3+3)^{-1}]$、
$H=\operatorname{Spec}\mathbb C[h,(h^6+27)^{-1}]$、
$T=\Lambda\times H$とする。
$F_\lambda=(x^3+y^3+z^3)/3+u^3+v^3+\lambda uxy$は
$\lambda^3+3\ne0$で滑らかなcubic threefoldを定める。
既存の混合族（\texttt{mt:eq:relations}, \texttt{mt:eq:center}）の
非可換代数$A_{h,\lambda}$と中心元$\Omega_{h,\lambda}$について、
$B_{h,\lambda}=A_{h,\lambda}/(\Omega_{h,\lambda})$と置く。

可換化した中心元は、$\mathbb P^4_T$上の三次式
\begin{equation}
 \Omega_{h,\lambda}^{\mathrm{ab}}
 =F_\lambda-\frac{2h^2}{3}xyz
 +\frac{h^2\lambda^2u^2}{6}z+\frac{h^4\lambda u}{18}z^2
 \tag{PL1}\label{psl:eq:abelian}
\end{equation}
である。次の閉部分スキームを定める：
\begin{align}
 \mathcal E&=V(u,v,x^3+y^3+z^3+h^2xyz),\notag\\
 \mathcal S&=V(x^2+\lambda uy,\ y^2+\lambda ux,\ z^2,
                       \Omega_{h,\lambda}^{\mathrm{ab}})
                  \subset\mathbb P^4_T.
 \tag{PL2}\label{psl:eq:families}
\end{align}

\begin{lemma}[孤立部分の24次の自由模型]\label{psl:lem:free}
$\mathcal S$は$u\ne0$に含まれる。$u=1$では
\begin{equation}
 \begin{split}
 \mathcal O(\mathcal S)
 &=R[x,y,z,v]/(x^2+\lambda y,\ y^2+\lambda x,\ z^2,\
                 1+v^3+\lambda xy/3+h^2zC),\\
 C&=\lambda^2/6-2xy/3,
 \quad R=\mathbb C[\lambda,h,(\lambda^3+3)^{-1},(h^6+27)^{-1}].
 \end{split}
 \tag{PL3}\label{psl:eq:24}
\end{equation}
この環は$R$上階数24の自由加群であり、基底は
\[
 \{x^ay^bz^cv^d\mid 0\le a,b,c\le1,\ 0\le d\le2\}.
 \tag{PL4}\label{psl:eq:basis}
\]
さらに$\mathcal S$は$h$に依らず抽象スキームとして自明化できる。
具体的に$u=1$上で$v$は単元であり、
\begin{equation}
 v'=v+\frac{h^2z}{3v^2}\left(\frac{\lambda^2}{6}-\frac{2xy}{3}\right)
 \tag{PL5}\label{psl:eq:shear}
\end{equation}
により$\mathcal S\simeq\mathcal S_{h=0}\times_\Lambda T$を得る。
\end{lemma}
\begin{proof}
$u=0$上で最初の三つの二次方程式は$x^2=y^2=z^2=0$となる。
その商では$\Omega^{\mathrm{ab}}=v^3$に冪零元を加えたものなので、
射影点では$v\ne0$の局所化と両立せず、$\mathcal S$の台は$u\ne0$である。
これは有限型でのset-theoreticな記述に加えて、$u=0$の閉部分スキームが
空であることを示す。

$u=1$では、$x^2=-\lambda y$、$y^2=-\lambda x$ゆえ
$x^3=y^3=-\lambda xy$、$z^2=z^3=0$となり、
\eqref{psl:eq:abelian}は\eqref{psl:eq:24}に一致する。
二つの方程式$x^2+\lambda y,y^2+\lambda x$は、総次数順序で
先頭単項式$x^2,y^2$を持つmonic Gr\"obner基底である。
実際唯一の$S$-多項式は$\lambda(y^3-x^3)$であり、
二つの関係により零へ還元される。
従って$R[x,y]$の商は基底$1,x,y,xy$を持つ階数4の自由加群。
その上に$z^2=0$を加えて階数8、最後に$v$のmonic三次関係を
加えて階数24を得る。従って$\mathcal S/T$は有限平坦である。

$t=xy$とすると$t^2=\lambda^2t$なので
\[
 (1+\lambda t/3)^{-1}=1-\frac{\lambda t}{3+\lambda^3}.
\]
従って$1+\lambda t/3$は単元、$h^2zC$は二乗零である。
式\eqref{psl:eq:24}から$v^3$、従って$v$も単元になる。
さらに$z^2=0$なので
\[
 (v')^3=v^3+h^2zC,
\]
となり、残りの関係は$1+(v')^3+\lambda xy/3=0$である。
逆向きの変換は同じ二乗零の変位を引けばよい。
これは$R$上の環の同型であり、単なる点集合の対応ではない。
\end{proof}

\begin{theorem}[任意の滑らかな混合係数での平坦極限]
\label{psl:thm:universal-flat}
$T$上の$\mathcal E$、$\mathcal S$は交わらず、
$\mathcal Z=\mathcal E\sqcup\mathcal S\subset\mathbb P^4_T$は平坦である。
すべての$\lambda^3+3\ne0$、$h\ne0$、$h^6+27\ne0$に対し
\[
 \mathcal Z_{h,\lambda}\simeq
 \operatorname{Point}(B_{h,\lambda})
 \qquad\text{（初期点の表現スキーム）},
 \tag{PL6}\label{psl:eq:point}
\]
$h=0$では$\mathbb P^4$内の閉部分スキームとして
\[
 \boxed{\quad \mathcal Z_{0,\lambda}=D_0(\pi_{X_\lambda}).\quad}
 \tag{PL7}\label{psl:eq:limit}
\]
従って一般$h$のpoint schemeのスキーム論的平坦閉包は
任意の滑らかな混合係数$\lambda$でPoisson零点スキームへ特殊化する。
全ファイバーのHilbert多項式は$3t+24$である。
\end{theorem}
\begin{proof}
補題\ref{psl:lem:free}より$\mathcal S$は$u\ne0$に含まれ、
$\mathcal E$は$u=0$上にあるから両者は交わらない。
$\mathcal S$は有限自由、$\mathcal E$はHesse pencilの相対平面三次曲線で
$h^6+27$を反転した範囲では滑らか、従って両者は$T$上平坦である。

$h\ne0$では中心一次元$u,v$の多重線形化により、どちらかが非零な
point moduleは定数列となる。$u=1$上のpoint moduleは
$B_{h,\lambda}$の可換化とスキームとして一致するので$\mathcal S$で表される。
$u=0,v\ne0$のチャートでは可換化の解はない。
$u=v=0$では$B_{h,\lambda}$は$\lambda$に依らない
三次元Sklyanin中心商となり、そのpoint schemeは滑らかな$\mathcal E_h$。
$h\ne0$ではSklyaninの点列shiftは非自明なtranslationだから
$p_0\ne p_1$となる。この二点の2次小行列は各点近傍で単元であり、
中心性の多重線形化から$u=v=0$が冪零的方向も含めて強制される。
従って$\mathcal E$に余分な厚化はない。
ここでは前節のSklyanin定理を使い、全$\lambda$に同じ証明を適用している。

$h=0$のPoisson bracketは
$\{y,z\}=x^2+\lambda uy$、
$\{z,x\}=y^2+\lambda ux$、
$\{x,y\}=z^2$、$\{u,-\}=\{v,-\}=0$。
$u=v=0$では$F_\lambda|_{u=v=0}$が滑らかなFermat三次曲線で、
EulerベクトルがアフィンPoissonの像に入るため射影化で零になる。
中心方向が非零の点ではEulerベクトルは像に入らず、
零点のスキーム条件は三つのactive偏微分と$F_\lambda$で与えられる。
これらはまさに$\mathcal S_{h=0}$の方程式である。
よって$\mathcal Z_{0,\lambda}=D_0(\pi_{X_\lambda})$は埋め込みスキームとして成立する。

$T$は整域、$\mathcal Z/T$は平坦なので$h$-torsionを持たず、
$D(h)$上の閉部分スキームのschematic closureに一致する。
曲線$\mathcal E_h$はHilbert多項式$3t$、有限部分は長さ24だから、
全体は$3t+24$である。
\end{proof}

\subsection{零混合係数での衝突}

\begin{corollary}[12個の二重点から3個の八重点への合体]
\label{psl:cor:collision}
$h\ne0$、$h^6+27\ne0$を固定する。
$\lambda\ne0$、$\lambda^3+3\ne0$では$\mathcal S_{h,\lambda}$は
12個の相異なる支持点を持ち、各局所環は
$\mathbb C[\varepsilon]/(\varepsilon^2)$に同型である。
$\lambda=0$では支持は三点となり、各局所環は
\[
 \boxed{\mathbb C[x,y,z]/(x^2,y^2,z^2)}
 \tag{PL8}\label{psl:eq:octuple}
\]
に同型で長さ8である。この衝突は有限平坦族の中で起き、
全長24が保存される。
\end{corollary}
\begin{proof}
$\lambda\ne0$では$x^2+\lambda y=y^2+\lambda x=0$の解は
$(x,y)=(0,0)$と、$x^3=-\lambda^3$、$y=-x^2/\lambda$を満たす三点である。
Jacobian行列式$4xy-\lambda^2$はそれぞれ
$-\lambda^2$、$3\lambda^2$で零でない。
$v'^3=-1-\lambda xy/3$は$xy=0$または$\lambda^2$なので、
$\lambda^3+3\ne0$なら各$x,y$解に三つの単純根が対応する。
$z^2=0$により各点の局所環はdual number。
$\lambda=0$では$x^2=y^2=z^2=0$、$v'^3=-1$となるので、
三つの根ごとに\eqref{psl:eq:octuple}を得る。
\end{proof}

\begin{remark}[特殊化と新規性の範囲]
\label{psl:rem:scope}
$\lambda=0$でもPoisson零点スキームの平坦極限は成立するが、
$B_{0,\lambda}$自身のpoint schemeは可換cubic $X_\lambda$全体である。
$\lambda^3=-3$では$X_\lambda$が特異となるため、上の滑らかな
Poisson cubicに関する主張の対象外である。
Sklyaninの退化値$h^6=-27$も除外した。
一般の非可換cubicがこのような平坦極限を持つとは主張しない。
上の全$\lambda$に関する定理は、特定係数での比較を
パラメータを動かす有限平坦な族へ拡張する。
この拡張の先行研究上の初出は未確認である。
\end{remark}

% 2026-11-11 / v42
\section{混合係数の判別式上の孤立point scheme}
\label{pdb:sec}
混合量子化族の関係式は
\[
[y,z]=h(x^2+\lambda uy),\quad[z,x]=h(y^2+\lambda ux),
\quad[x,y]=hz^2,\quad [u,-]=[v,-]=0.
\]
中心三次元の可換化は
\[
\Omega_{h,\lambda}^{\rm ab}
=F_\lambda-\tfrac{2h^2}{3}xyz+\tfrac{h^2\lambda^2u^2}{6}z
 +\tfrac{h^4\lambda u}{18}z^2,\quad
F_\lambda=(x^3+y^3+z^3)/3+u^3+v^3+\lambda uxy.
\]
$h\ne0$で中心座標$u\ne0$のpoint moduleは定数列となるので、
$u=1$チャートの可換化の商環で表される。以下の環自体は$h=0$でも定義する。
\begin{equation}
\mathcal T=\mathbb C[\lambda,h,x,y,z,v]/
(x^2+\lambda y,\ y^2+\lambda x,\ z^2,\
1+v^3+\lambda xy/3+h^2z(\lambda^2/6-2xy/3)).
\label{pdb:eq:ring}
\end{equation}

\begin{theorem}[判別式を越える有限自由性]\label{pdb:thm:flat}
$\mathcal T$は$\mathbb C[\lambda,h]$上自由階数24で、基底は
\[
\{x^ay^bz^cv^d:0\le a,b,c\le1,\quad0\le d\le2\}.
\]
したがって$\lambda^3=-3$上でも有限平坦性は失われない。
\end{theorem}
\begin{proof}
$x^2+\lambda y,y^2+\lambda x$は総次数順序で
先頭単項式$x^2,y^2$を持つmonic Gr\"obner基底である。
唯一の$S$多項式$\lambda(y^3-x^3)$は
$x^3=y^3=-\lambda xy$により零へ還元される。
従って最初の二式の商は基底$1,x,y,xy$を持つ階数4の
$\mathbb C[\lambda,h]$自由加群。
$z^2$を加えて階数8、最後の$v$に関するmonic三次関係で
$1,v,v^2$を加えて階数24となる。この証明は分母を使わない。
\end{proof}

\begin{theorem}[三個の二重点と三個の六重点]\label{pdb:thm:sixfold}
$\lambda_0^3=-3$かつ$h_0\ne0$とする。
式\eqref{pdb:eq:ring}の$(\lambda,h)=(\lambda_0,h_0)$特殊化は
六つの支持点を持つ。
$x=y=0$上には$v^3=-1$の三個の長さ2の点があり、
各局所環は$\mathbb C[\varepsilon]/(\varepsilon^2)$である。
他の三点は
\[
 (x,y)=(-\lambda_0\zeta,-\lambda_0\zeta^2),\qquad \zeta^3=1
\]
に支持され、各局所環は$\mathbb C[t]/(t^6)$である。
従って長さは$3\cdot2+3\cdot6=24$。
同じ$\lambda_0$で$h=0$なら後者の局所環は
$\mathbb C[t,z]/(t^3,z^2)$となり、同型ではない。
\end{theorem}
\begin{proof}
$\lambda_0\ne0$なので最初の二式の解は
$(0,0)$または表示した三つの点である。
Jacobian行列式$4xy-\lambda_0^2$は前者で$-\lambda_0^2$、
後者で$3\lambda_0^2$だから四つの枝はetaleに分離する。
$x=y=0$では
$1+v^3+(h_0^2\lambda_0^2/6)z=0$、$z^2=0$であり、
$v^3=-1$の単純な三根それぞれで長さ2を得る。
他の枝では$xy=\lambda_0^2$、$\lambda_0^3=-3$なので
\[
 1+\lambda_0xy/3=0,\qquad
 \lambda_0^2/6-2xy/3=-\lambda_0^2/2.
\]
従って各局所環は
\[
 \mathbb C[v,z]/(z^2,v^3-cz),\qquad c=h_0^2\lambda_0^2/2.
\]
$c\ne0$なら$z=v^3/c$を消去し$\mathbb C[v]/(v^6)$、
$c=0$なら$\mathbb C[v,z]/(v^3,z^2)$。
両者のZariski接空間の次元はそれぞれ1,2である。
\end{proof}

\begin{corollary}[衝突の局所模型]\label{pdb:cor:double-cubic}
$\lambda^3+3=\delta$と置き、$\lambda\ne0$で
$(x,y)=(-\lambda\zeta,-\lambda\zeta^2)$の枝を固定する。
その有限部分の局所模型は
\[
 \mathbb C[\lambda,h,v,z]_{\rm loc}/
 (z^2,\ v^3+\delta/3-(h^2\lambda^2/2)z).
\]
$h\ne0$なら$z$を消去し
\[
\boxed{\mathbb C[\lambda,h,v]_{\rm loc}/((v^3+\delta/3)^2)}
\tag{DB}\label{pdb:eq:double-cubic}
\]
となる。すなわち$\delta\ne0$の三個の二重点が、
$\delta=0$で一個の六重点に合体する。
\end{corollary}
\begin{proof}
三つの非零枝では$xy=\lambda^2$なので、式
\eqref{pdb:eq:ring}の最後の方程式は
$v^3+\delta/3-(h^2\lambda^2/2)z=0$。
$h\lambda\ne0$なら
$z=2(v^3+\delta/3)/(h^2\lambda^2)$を代入し、$z^2=0$
から表示を得る。
\end{proof}

\begin{remark}[適用範囲]\label{pdb:rem:scope}
$\lambda^3=-3$では可換三次超曲面自身が特異なので、
滑らかなPoisson cubicの零点定理やBondal予想の
適用対象ではない。本節は中心座標非零チャートの
point moduleを記述する可換化環の平坦性と局所型を証明した。
$h\ne0$かつ$h^6\ne-27$では、別に得た
Sklyaninのpoint schemeの既知表示と併せて
全point scheme中の有限部分とみなせる。
$h=0$での環\eqref{pdb:eq:ring}は$B_{0,\lambda}$自身の
point scheme（可換cubic全体）とは異なる。
$h^6=-27$でのSklyanin成分の退化、任意のnc cubicへの拡張は未解決である。
\end{remark}

\begin{thebibliography}{99}
\bibitem{PSPSmithTate}
S.~P.~Smith and J.~Tate,
\emph{The center of the 3-dimensional and 4-dimensional Sklyanin algebras},
K-Theory \textbf{8} (1994), 19--63, Section~1 (pp.~20--21).
\href{https://sites.math.washington.edu/~smith/Research/tate.pdf}{著者公開版}.
本稿では中心三次元による楕円曲線のtwisted homogeneous coordinate ringへの商を用いる。
有限位数を仮定する同論文後半の中心の分類は用いない。

\bibitem{PSPDeLaet}
K.~De Laet,
\emph{On the center of 3-dimensional and 4-dimensional Sklyanin algebras},
J. Algebra \textbf{487} (2017), 244--268,
\href{https://doi.org/10.1016/j.jalgebra.2017.05.032}{doi:10.1016/j.jalgebra.2017.05.032}.
参照箇所は\href{https://arxiv.org/abs/1612.06158}{arXiv:1612.06158v1}の
Definition~4.1、Sections~4.1--4.2、Theorem~4.4。
点曲線、平行移動$\tau=[a:b:c]$、中心三次元による点加群の消去を用いる。

\end{thebibliography}
\end{document}
